\documentclass[aspectratio=169,10pt]{beamer}

% Theme
\usetheme{Madrid}
\usecolortheme{whale}
\usefonttheme{structurebold}
\setbeamertemplate{navigation symbols}{}
\setbeamertemplate{footline}[frame number]

% Packages
\usepackage[italian]{babel}
\usepackage[utf8]{inputenc}
\usepackage[T1]{fontenc}
\usepackage{lmodern}
\usepackage{amsmath,amssymb}
\usepackage{booktabs}
\usepackage{xcolor}
\usepackage{listings}
\usepackage{tikz}
\usetikzlibrary{shapes.geometric,arrows.meta,positioning,fit,backgrounds}
\usepackage{pgfplots}
\pgfplotsset{compat=newest}
\usepackage{graphicx}
\usepackage{multirow}
\usepackage{array}
\usepackage{colortbl}

% Colours
\definecolor{riverblue}{RGB}{31,119,180}
\definecolor{capyorange}{RGB}{214,118,27}
\definecolor{driftred}{RGB}{180,40,40}
\definecolor{adwingreen}{RGB}{44,160,44}
\definecolor{codebg}{RGB}{245,245,245}
\definecolor{javakw}{RGB}{0,0,200}
\definecolor{javastr}{RGB}{163,21,21}
\definecolor{javacomm}{RGB}{0,128,0}
\definecolor{pythonkw}{RGB}{0,0,200}

% Code style
\lstdefinestyle{java}{
  language=Java,
  basicstyle=\ttfamily\fontsize{6.5}{8}\selectfont,
  backgroundcolor=\color{codebg},
  keywordstyle=\color{javakw}\bfseries,
  stringstyle=\color{javastr},
  commentstyle=\color{javacomm}\itshape,
  frame=single,framerule=0.4pt,
  breaklines=true,
  tabsize=2,
  showstringspaces=false,
  xleftmargin=4pt, xrightmargin=4pt,
  numbers=none,
}
\lstdefinestyle{python}{
  language=Python,
  basicstyle=\ttfamily\fontsize{6.5}{8}\selectfont,
  backgroundcolor=\color{codebg},
  keywordstyle=\color{pythonkw}\bfseries,
  stringstyle=\color{javastr},
  commentstyle=\color{javacomm}\itshape,
  frame=single,framerule=0.4pt,
  breaklines=true,
  tabsize=2,
  showstringspaces=false,
  xleftmargin=4pt, xrightmargin=4pt,
  numbers=none,
}

% Macros
\newcommand{\hatr}{\textsc{hatr}}
\newcommand{\river}{\textcolor{riverblue}{\textbf{River}}}
\newcommand{\capymoa}{\textcolor{capyorange}{\textbf{CapyMOA}}}
\newcommand{\adwin}{\textsc{adwin}}
\newcommand{\fimtdd}{\textsc{fimtdd}}
\newcommand{\arf}{\textsc{arf}}
\newcommand{\eps}{\varepsilon}

% Title
\title[\hatr{} — Porting e Validazione]{%
  Hoeffding Adaptive Tree Regressor\\
  Teoria, Implementazione in MOA/CapyMOA e Validazione Sperimentale}
\subtitle{Stream Data Analytics 2025--2026}
\author[Andrea Zhang]{Andrea Zhang}
\institute[Polimi]{Politecnico di Milano}
\date{2026}

\begin{document}

\begin{frame}
  \titlepage
\end{frame}

\begin{frame}{Indice}
  \tableofcontents[hideallsubsections]
\end{frame}

% ============================================================
\section{Motivazione}
% ============================================================

\begin{frame}{Perché portare \hatr{} in CapyMOA?}
  \begin{columns}[T]
    \column{0.45\linewidth}
    \begin{alertblock}{Il gap}
      \small
      \river{} offre un'implementazione di riferimento di \hatr{},
      ma MOA/CapyMOA non dispone di un learner nativo equivalente.

      \smallskip
      Il confronto richiede quindi di cambiare framework, evaluator
      e pipeline sperimentale.
    \end{alertblock}

    \vspace{-2pt}
    \begin{center}
      \Large\textcolor{riverblue}{$\Downarrow$}
    \end{center}
    \vspace{-6pt}

    \begin{block}{Il contributo}
      \small
      Un port Java nativo con:
      \begin{itemize}
        \item wrapper Python compatibile con CapyMOA;
        \item validazione rispetto a River su 8 dataset;
        \item integrazione nel ciclo di vita MOA.
      \end{itemize}
    \end{block}

    \column{0.51\linewidth}
    \begin{block}{Prestazioni osservate}
      \small
      \begin{itemize}
        \item \textbf{Training}: circa 3--8$\times$ più veloce nel benchmark;
        \item \textbf{Memoria}: modello generalmente più compatto.
      \end{itemize}
    \end{block}

    \begin{block}{Integrazione nel framework}
      \small
      \begin{itemize}
        \item \textbf{Schema ARFF}: feature nominali riconosciute
              automaticamente;
        \item \textbf{Ecosistema MOA}: stessi stream, evaluator e learner
              di confronto.
      \end{itemize}
    \end{block}

    \vspace{1pt}
    {\scriptsize\textcolor{gray}{Il vantaggio riguarda soprattutto il training;
    nell'inferenza per singola istanza può pesare l'overhead Python--JVM.}}
  \end{columns}
\end{frame}

% ============================================================
\section{Background Teorico}
% ============================================================

\begin{frame}{Hoeffding Tree: base per lo streaming}
  \begin{columns}[T]
    \column{0.55\linewidth}
    \textbf{Hoeffding Tree (VFDT)} \cite{Domingos2000} è un albero decisionale
    online che costruisce split usando la \alert{Hoeffding bound}:
    \[
      \varepsilon(n, \delta) = \sqrt{\frac{\ln(1/\delta)}{2n}}
    \]
    \medskip
    \textbf{Regola di split}: dati i due migliori attributi $a_1, a_2$
    (per variance reduction):
    \[
      \text{split se}\quad
      \frac{\Delta\bar{G}(a_2)}{\Delta\bar{G}(a_1)} < 1 - \varepsilon
      \quad\text{oppure}\quad \varepsilon < \tau
    \]
    \medskip
    \textbf{Grace period} $g$: si tenta lo split ogni $g$ istanze per
    ammortizzare il costo.

    \column{0.42\linewidth}
    \begin{block}{Split criterion — Variance Reduction}
      \vspace{2pt}
      \[
        \Delta\bar{G}(a) = \mathrm{Var}(S) - \sum_{v} \frac{|S_v|}{|S|} \mathrm{Var}(S_v)
      \]
      \vspace{-4pt}
      \begin{itemize}\small
        \item $S$ = istanze nel nodo
        \item $S_v$ = sottoinsieme che va al ramo $v$
        \item $\mathrm{Var}$ = varianza campionaria (\emph{ddof}=1, Welford)
      \end{itemize}
    \end{block}
    \begin{alertblock}{Limitazione}
      Un Hoeffding Tree standard non rileva i \textbf{concept drift}:
      una volta creato, un nodo non viene revisionato.
    \end{alertblock}
  \end{columns}
\end{frame}

\begin{frame}{ADWIN: finestra adattiva per il drift}
  \begin{columns}[T]
    \column{0.52\linewidth}
    \textbf{ADWIN} (ADaptive WINdowing) \cite{Bifet2007} mantiene una
    \alert{finestra di dimensione variabile} sullo stream di errori.

    \medskip
    \textbf{Idea}: si cerca il taglio $W = W_0 \cup W_1$ che massimizza
    la differenza di medie; se questa supera la soglia, la parte vecchia viene
    eliminata.

    \medskip
    \textbf{Criterio di taglio} (Babcock et al., 2003 — variance-aware):
    \[
      |\hat{\mu}_{W_0} - \hat{\mu}_{W_1}| > \varepsilon_{cut}
    \]
    \[
      \varepsilon_{cut} = \sqrt{2\,m^{-1}\,\hat{\sigma}^2_W\,\delta'}
                        + \frac{2}{3}\,m^{-1}\,\delta'
    \]
    \[
      m^{-1} = \frac{1}{n_0-n_{\min}+1}
             + \frac{1}{n_1-n_{\min}+1}
    \]
    dove $n_0=|W_0|$, $n_1=|W_1|$,
    $n_{\min}=\texttt{minWindowLength}$ e
    $\delta' = \ln\!\bigl(2\ln(|W|)/\delta\bigr)$.

    \column{0.45\linewidth}
    \begin{block}{Struttura a bucket}
\centering
\begin{tikzpicture}[scale=0.95]

\tikzset{
  bucket/.style={
    rectangle,
    draw,
    minimum width=0.75cm,
    minimum height=0.45cm,
    fill=blue!10,
    font=\scriptsize
  }
}

% Row 0
\node[left,font=\scriptsize] at (-0.8,1.6) {row 0};
\node[bucket] (b0a) at (0,1.6) {$2^0$};
\node[bucket] (b0b) at (0.9,1.6) {$2^0$};
\node[bucket] (b0c) at (1.8,1.6) {$2^0$};

% Row 1
\node[left,font=\scriptsize] at (-0.8,0.8) {row 1};
\node[bucket] (b1a) at (0,0.8) {$2^1$};
\node[bucket] (b1b) at (0.9,0.8) {$2^1$};

% Row 2
\node[left,font=\scriptsize] at (-0.8,0.0) {row 2};
\node[bucket] (b2a) at (0,0.0) {$2^2$};

% Compress arrow
\draw[->,thick,driftred]
  (2.8,1.6) -- (2.8,0.8)
  node[midway,right,font=\scriptsize]{compress};

% Direction label
\node[font=\small,text=driftred]
  at (1.0,-0.8)
  {$\leftarrow$ newest \qquad oldest $\rightarrow$};

\end{tikzpicture}

\vspace{2mm}

Ogni riga $i$ contiene bucket di dimensione $2^i$; quando una riga è
piena, il \emph{compress} fonde i due bucket più vecchi nella riga
successiva.

\end{block}
  \end{columns}
\end{frame}

\begin{frame}{Meccanismo adattivo di \hatr{}}
  \begin{columns}[T]
    \column{0.48\linewidth}
    Ogni nodo monitora l'errore con \adwin{}. Un \textbf{nodo interno}
    può inoltre ospitare:
    \begin{enumerate}
      \item lo storico dell'errore $|y - \hat{y}|$
      \item un \textbf{albero alternativo} opzionale, inizialmente assente
    \end{enumerate}

    \medskip
    \textbf{Ciclo di vita} a ogni istanza $(\mathbf{x}, y)$:
    \begin{itemize}
      \item Predici dal nodo corrente; calcola errore
      \item Aggiorna \adwin{}
      \item \alert{Se drift rilevato} e l'errore \emph{sta aumentando}:
            crea albero alternativo (foglia) alla stessa profondità
      \item \alert{Se entrambi i path sono maturi}
            ($n_a,n_c > w_{\min}$): esegui \textbf{z-test}
      \item Se z-test significativo: \emph{swap} o \emph{pruning}
    \end{itemize}

    \column{0.49\linewidth}
    \begin{alertblock}{Z-test per lo swap}
      \vspace{2pt}
      Confronta l'errore medio del path corrente ($\mu_c, \sigma_c^2$)
      con quello dell'alternativo ($\mu_a, \sigma_a^2$):
      \[
        z = \frac{\mu_a - \mu_c}{\sqrt{\sigma_a^2/n_a + \sigma_c^2/n_c}}
      \]
      \[
        p = 2\,\Phi(-|z|)
      \]
      \begin{itemize}\small
        \item $p \le \alpha$ e $\mu_a < \mu_c$: \textbf{swap} (alternativo è migliore)
        \item $p \le \alpha$ e $\mu_a \ge \mu_c$: \textbf{pruning} (corrente è migliore)
      \end{itemize}
    \end{alertblock}

    \vspace{2pt}
    \textbf{Predizione}: finché esistono alternative, media le predizioni
    delle foglie raggiunte lungo tutti i path attivi.
  \end{columns}
\end{frame}

\begin{frame}{Modalità di foglia}
  \hatr{} supporta tre strategie di predizione alle foglie:

  \bigskip
  \begin{columns}[T]
    \column{0.3\linewidth}
    \begin{block}{\textbf{MEAN}}
      Predice con la \alert{media target} accumulata nella foglia
      (aggiornamento incrementale).
      \[
        \hat{y} = \bar{y}_{\text{leaf}}
      \]
      \vspace{2pt}\\
      \textit{Semplice, robusto, deterministico.}
    \end{block}

    \column{0.33\linewidth}
    \begin{block}{\textbf{MODEL}}
      Predice con un \alert{modello lineare online} (\emph{Perceptron},
      gradiente su perdita quadratica):
      \[
        \hat{y} = \mathbf{w}^{\!\top}\!\mathbf{x} + b
      \]
      \[
        \nabla = 2(\hat{y}-y)\,[\mathbf{x};\,1]
      \]
      \vspace{2pt}\\
      I figli ereditano una \textbf{copia} del modello del padre.
    \end{block}

    \column{0.33\linewidth}
    \begin{block}{\textbf{ADAPTIVE} (default)}
      Sceglie dinamicamente tra \textsc{mean} e \textsc{model} tramite
      \alert{faded MSE} (EWMA):
      \[
        e_{\text{mean}} \leftarrow \gamma\,e_{\text{mean}} + (y-\bar{y})^2
      \]
      \[
        e_{\text{model}} \leftarrow \gamma\,e_{\text{model}} + (y-\hat{y}_m)^2
      \]
      Usa \textsc{mean} se $e_{\text{mean}} < e_{\text{model}}$.
    \end{block}
  \end{columns}

  \bigskip
  \begin{center}
    \small
    Bootstrap sampling opzionale (Poisson($\lambda=1$)) in tutte le
    modalità per aumentare la diversità delle foglie.
  \end{center}
\end{frame}

% ============================================================
\section{Scelte di Design e Risultati}
% ============================================================

\begin{frame}[shrink=15]{Scelte di design: configurazione implementata}
  \begin{columns}[T]
    \column{0.50\linewidth}
    \begin{block}{Cosa è stato implementato}
      \small
      \textbf{Solo la configurazione default di River}:
      \begin{itemize}
        \item Splitter: \texttt{HATEBSTObserver} (Truncated E-BST)\\
              {\footnotesize split binari per attributi numerici}
        \item Drift detector: \texttt{ADWINDetector}
        \item Leaf model: \texttt{HAPerceptron} (SGD lineare)
        \item \texttt{binary\_split=False} hardcoded\\
              {\footnotesize attributi nominali $\to$ multiway branch}
      \end{itemize}
      \smallskip
      Le principali opzioni algoritmiche sono configurabili
      dal wrapper Python e tradotte in flag MOA.
    \end{block}

    \medskip
    \begin{block}{Motivazione}
      \small
      Il port replica la configurazione standard di \river{}.
      Senza bootstrap produce le stesse predizioni; con bootstrap
      i diversi generatori casuali rendono confrontabili le metriche,
      non le singole predizioni.
    \end{block}

    \column{0.47\linewidth}
    \begin{alertblock}{Cosa non è stato implementato}
      \small
      \begin{itemize}
        \item \textbf{\texttt{binary\_split=True}}\\
              {\footnotesize split nominale binario forzato
              (non-default in River)}
        \item \textbf{\texttt{QOSplitter}}\\
              {\footnotesize richiederebbe
              \texttt{AdaNumMultiwayBranch} per attributi numerici}
        \item \textbf{Componenti pluggabili}\\
              {\footnotesize splitter, drift detector e leaf model
              sono fissi al default River}
      \end{itemize}
    \end{alertblock}

    \medskip
    \begin{tabular}{@{}ll@{\;}c@{}}
      \toprule
      \textbf{Stato} & \textbf{Esempio} & \\
      \midrule
      Configurabile & \texttt{grace\_period}, \texttt{delta} & \textcolor{adwingreen}{$\checkmark$} \\
      Hardcoded default & \texttt{splitter}, \texttt{ADWIN} & \textcolor{capyorange}{$\sim$} \\
      Non implementato & \texttt{binary\_split=True} & \textcolor{driftred}{$\times$} \\
      \bottomrule
    \end{tabular}
  \end{columns}
\end{frame}

\begin{frame}{Setup sperimentale}
  \begin{columns}[T]
    \column{0.48\linewidth}
    \textbf{Dataset} ($n=17\,000$, prequential evaluation):
    \smallskip
    \footnotesize
    \begin{tabular}{@{}lcc@{}}
      \toprule
      Dataset & Feat. & Drift \\
      \midrule
      Synthetic    & 5  & abrupt \\
      Bike Sharing & 12 & naturale \\
      Fried        & 10 & nessuno \\
      Hyper (lento)& 10 & graduale \\
      HyperFast    & 10 & rapido \\
      FriedDrift   & 10 & abrupt \\
      RBFReg       & 10 & graduale \\
      ElecDemand   & 7  & naturale \\
      \bottomrule
    \end{tabular}

    \smallskip
    {\footnotesize\textcolor{gray}{Cap a 17\,000 ist.
    Bike/ElecDemand: reali. Hyper/HyperFast: generatore.}}

    \smallskip
    \footnotesize\textbf{Configurazioni}:
    \begin{itemize}
      \item \textit{Deterministic}: \texttt{mean}, no bootstrap
      \item \textit{Default}: \texttt{adaptive}, bootstrap=True
    \end{itemize}

    \column{0.49\linewidth}
    \vspace{-2pt}
    \textbf{Iperparametri comuni}:
    \begin{itemize}\small
      \item \texttt{grace\_period=50}, $\delta=10^{-7}$, $\tau=0.05$
      \item \texttt{adwin\_delta=0.002}, \texttt{switch\_significance=0.05}
      \item \texttt{drift\_window\_threshold=300}
      \item \texttt{model\_selector\_decay=0.95}
      \item Leaf model: $lr_w=lr_b=0.01$, $L2=0$ (default fisso); seed=0
      \item StandardScaler online identico ai due lati
    \end{itemize}

    \vspace{4pt}
    \begin{alertblock}{Due domande diverse}
      \small
      La configurazione deterministica misura la \textbf{fedeltà del port};
      quella default misura la \textbf{qualità d'uso reale}.
    \end{alertblock}
  \end{columns}
\end{frame}

\begin{frame}{Quattro risultati da verificare nei notebook}
  \footnotesize
  \begin{tabular}{@{}p{0.28\linewidth}p{0.30\linewidth}p{0.36\linewidth}@{}}
    \toprule
    \textbf{Domanda} & \textbf{Esperimento} & \textbf{Risultato} \\
    \midrule
    Il port replica River? &
      MEAN/MODEL/ADAPTIVE senza bootstrap &
      Predizioni e struttura coincidono numericamente \\
    \addlinespace
    Il default mantiene la qualità? &
      ADAPTIVE con bootstrap &
      MAE e RMSE restano comparabili; le predizioni divergono per la RNG \\
    \addlinespace
    Il training è più rapido? &
      8 dataset, cap a 17\,000 istanze &
      Speedup complessivo osservato: circa 2.95--7.85$\times$ \\
    \addlinespace
    Il modello è più compatto? &
      Stima finale del modello &
      Riduzione osservata: circa 2.2--11.9$\times$\textsuperscript{*} \\
    \bottomrule
  \end{tabular}

  \medskip
  {\footnotesize\textsuperscript{*}\,Confronto indicativo: gli strumenti
  misurano heap Python e oggetti JVM con metodologie differenti.}

  \medskip
  \begin{alertblock}{Demo}
    \footnotesize
    I notebook mostrano dati, configurazioni, metriche e curve complete;
    questa tabella è la mappa con cui leggere la dimostrazione.
  \end{alertblock}
\end{frame}

% ============================================================
\section{Dettagli Implementativi}
% ============================================================

\begin{frame}[fragile]{Struttura Java: gerarchia dei nodi}
  \centering
    \begin{tikzpicture}[
      cls/.style={rectangle,draw,fill=blue!8,font=\tiny\ttfamily,
                  text width=2.5cm,minimum height=0.5cm,align=center},
      abs/.style={cls,fill=gray!12,font=\tiny\ttfamily\itshape},
      arr/.style={-{Triangle[open,length=5pt,width=7pt]},thick}]
      \node[abs] (hanode) {HANode (interface)};
      \node[abs] (haleaf)   [below left=0.5cm and 0.3cm of hanode]  {HALeafNode (abstract)};
      \node[abs] (habranch) [below right=0.5cm and 0.1cm of hanode] {HABranchNode (abstract)};
      \node[abs] (adaleaf)  [below=0.5cm of haleaf]   {AdaLeafNode (abstract)};
      \node[abs] (adabranch)[below=0.5cm of habranch] {AdaBranchNode (abstract)};
      \node[cls] (alm)  [below left=0.5cm and -0.2cm of adaleaf]  {AdaLeafMean};
      \node[cls] (almo) [below=0.45cm of adaleaf]                  {AdaLeafModel};
      \node[cls] (alad) [below right=0.5cm and -0.15cm of adaleaf] {AdaLeafAdaptive};
      \node[cls,text width=1.8cm] (anb)
            [below=1.3cm of adabranch, xshift=-2.0cm] {AdaNum\\BinaryBranch};
      \node[cls,text width=1.8cm] (anombb)
            [below=1.3cm of adabranch]                {AdaNom\\BinaryBranch};
      \node[cls,text width=1.8cm] (anomm)
            [below=1.3cm of adabranch, xshift=2.0cm]  {AdaNom\\MultiwayBranch};
      \draw[arr] (haleaf) -- (hanode);
      \draw[arr] (habranch) -- (hanode);
      \draw[arr] (adaleaf) -- (haleaf);
      \draw[arr] (adabranch) -- (habranch);
      \draw[arr] (alm) -- (adaleaf);
      \draw[arr] (almo) -- (adaleaf);
      \draw[arr] (alad) -- (adaleaf);
      \draw[arr] (anb) -- (adabranch);
      \draw[arr] (anombb) -- (adabranch);
      \draw[arr] (anomm) -- (adabranch);
    \end{tikzpicture}

  \medskip
  \small
  \textbf{Separazione delle responsabilità:}
  le foglie apprendono e propongono split; i branch instradano le istanze
  e gestiscono gli alberi alternativi.
\end{frame}

\begin{frame}[fragile,shrink=6]{Codice: apprendimento in una foglia adattiva}
  \begin{columns}[T]
    \column{0.64\linewidth}
\begin{lstlisting}[style=java,basicstyle=\ttfamily\fontsize{7.1}{8.4}\selectfont]
double y = inst.classValue();
double w = inst.weight();
double yPred = getPrediction(inst);       // pre-update

if (tree.bootstrapSampling) {
  int k = poissonSample(rng);             // Poisson(1)
  if (k > 0) w *= k;
}

double err = Math.abs(y - yPred);
driftDetector.update(err);                // local ADWIN
errorTracker.update(err, 1.0);

updateModelSelector(inst, y, tree);       // faded MSE
baseLearnOne(inst, y, w);                 // stats + observer
trainLeafModel(inst, y, w);               // SGD model

double seen = getTotalWeight();
if (seen - lastSplitAttemptAt >= tree.gracePeriod
    && isActive() && tree.growthAllowed) {
  tree.attemptToSplit(this, parent, parentBranch,
                      tree.driftDetectorProto.createNew());
  lastSplitAttemptAt = seen;
}
\end{lstlisting}

    \column{0.33\linewidth}
    \begin{alertblock}{Perché l'ordine conta}
      \small
      Predizione, errore e selezione del modello usano lo stato
      \textbf{prima} dell'aggiornamento: così non valutano il target
      appena osservato con un modello che lo ha già appreso.
    \end{alertblock}

    \medskip
    \begin{block}{Grace period}
      \small
      Gli observer vengono aggiornati a ogni istanza, ma la ricerca
      dello split parte soltanto ogni $g$ unità di peso osservato.
    \end{block}
  \end{columns}
\end{frame}

\begin{frame}[fragile,shrink=5]{Codice: dalla Hoeffding bound al nuovo branch}
  \begin{columns}[T]
    \column{0.59\linewidth}
\begin{lstlisting}[style=java,basicstyle=\ttfamily\fontsize{7.1}{8.4}\selectfont]
Collections.sort(candidates);
SplitCandidate best = candidates.get(size - 1);
SplitCandidate second = candidates.get(size - 2);

double n = leaf.getTotalWeight();
double eps = Math.sqrt(
    Math.log(1.0 / delta) / (2.0 * n));

boolean shouldSplit = best.merit > 0.0 && (
    second.merit / best.merit < 1.0 - eps
    || eps < tau);

if (best.isNullSplit()) {
  leaf.deactivate();                       // pre-pruning
  return;
}

if (best.isNumeric)
  branch = new AdaNumBinaryBranch(...);
else if (best.isMultiway())
  branch = new AdaNomMultiwayBranch(...);
else
  branch = new AdaNomBinaryBranch(...);
\end{lstlisting}

    \column{0.37\linewidth}
    \begin{block}{Due condizioni di split}
      \small
      \begin{itemize}
        \item il rapporto tra secondo e primo candidato è abbastanza piccolo;
        \item oppure $\varepsilon < \tau$, per rompere una situazione di parità.
      \end{itemize}
    \end{block}

    \begin{alertblock}{Lo stesso risultato, tre strutture}
      \small
      Il tipo di \texttt{SplitCandidate} determina quale branch adattivo
      sostituisce la foglia, mantenendo statistiche e profondità.
    \end{alertblock}
  \end{columns}
\end{frame}

\begin{frame}{Tipi di branch (split)}

\vspace{-0.5em}

\begin{block}{Tipi di branch (split)}
  \begin{itemize}\small
    \item \textbf{AdaNumBinaryBranch}\\
          \footnotesize split numerico: $x_i \le t$;
          soglia da E-BST / TE-BST
    \medskip

    \item \textbf{AdaNomBinaryBranch}\\
          \footnotesize split nominale binario:
          $x_i = v$ vs.\ tutto il resto
    \medskip

    \item \textbf{AdaNomMultiwayBranch}\\
          \footnotesize un ramo per valore categorico;
          valore non visto $\to$ nuovo figlio dedicato
          (come River \texttt{add\_child});
          \texttt{maxBranches() = -1}
  \end{itemize}
\end{block}

\vspace{-0.8em}

\medskip
\begin{alertblock}{Perché non \texttt{AdaNumMultiwayBranch}?}
  \footnotesize
  In River esiste \texttt{AdaNumMultiwayBranchReg}, ma viene usata
  \emph{solo} con \texttt{QOSplitter(allow\_multiway\_splits=True)},
  configurazione non-default. Il default splitter (\texttt{TEBSTSplitter})
  produce sempre split binari per attributi numerici
  $\Rightarrow$ \texttt{AdaNumBinaryBranch} è sufficiente per
  l'equivalenza con River HATR standard.
\end{alertblock}

\vspace{-0.8em}

\medskip
\begin{alertblock}{\texttt{binary\_split} non implementato come opzione}
  \footnotesize
  River espone \texttt{binary\_split=False} (default): per attributi
  nominali il multiway viene preferito; il binario lo sostituisce solo
  se ha merit strettamente superiore.
  Il Java replica questo comportamento in modo fisso in
  \texttt{HANominalObserver}: \texttt{binary\_split=True} non è
  supportato.
\end{alertblock}

\end{frame}

\begin{frame}{Componenti di supporto}

\begin{block}{Observer per attributi numerici}
\begin{itemize}\small
  \item \texttt{HAEBSTObserver} — E-BST standard

  \item \texttt{HATEBSTObserver} — \textbf{Truncated} E-BST con arrotondamento
        \emph{half-even} (\texttt{BigDecimal.HALF\_EVEN}) per replicare
        \texttt{round(x, digits)} di Python
\end{itemize}
\end{block}

\vspace{0.2cm}

\begin{block}{Utilities auto-contenute}
\begin{itemize}\small
  \item \texttt{VarStats} — varianza incrementale Welford (port di
        \texttt{river.stats.Var})

  \item \texttt{ADWINDetector} — port fedele da River Cython

  \item \texttt{HAPerceptron} — modello lineare SGD

  \item \texttt{NormalDist} — CDF via Apache Commons Math 3
        (\texttt{Erf.erfc})
\end{itemize}
\end{block}

\vspace{0.2cm}

\end{frame}

\begin{frame}{Dettagli implementativi non ovvi}
  \begin{block}{Truncated EBST e arrotondamento}
    \small
    Il TEBST arrotonda i valori di split a un numero fisso di cifre
    decimali, riducendo i candidati unici nell'albero.
    Python \texttt{round()} usa banker's rounding (half-even);
    \texttt{Math.round()} in Java usa half-up.
    La discrepanza rompeva l'equivalenza: risolto usando
    \texttt{BigDecimal(x).setScale(d, HALF\_EVEN)}.
  \end{block}

  \medskip

  \begin{block}{Predizione con albero alternativo attivo}
    \small
    \texttt{traverseWithAlternate} non viene chiamato solo al momento
    dello swap: durante \emph{tutta la vita} dell'albero alternativo,
    la predizione è la media tra il path principale e quello alternativo.
  \end{block}
\end{frame}

\begin{frame}{Attributi nominali: River vs.\ MOA}
  \begin{block}{River — schema-less}
    \small
    Le feature arrivano come dizionari Python senza tipo dichiarato.
    \texttt{nominal\_attributes} (lista di nomi) è l'unico modo per
    indicare al tree quali feature trattare come categoriche; senza,
    verrebbero passate al \texttt{TEBSTSplitter} numerico.
  \end{block}

  \medskip

  \begin{block}{MOA — schema ARFF}
    \small
    Ogni dataset porta uno schema che dichiara il tipo di ogni attributo.
    \texttt{inst.attribute(i).isNominal()} funziona in automatico:
    nessuna configurazione richiesta all'utente.
  \end{block}
\end{frame}

\begin{frame}{\texttt{merit\_preprune}: un caso limite utile}

  \begin{block}{Meccanismo}
    \small
    Con \texttt{merit\_preprune=True} il candidato \emph{null split}
    (merit~$= -\infty$) viene aggiunto alla lista.
    Se vince il test di Hoeffding la foglia viene \textbf{disattivata}.

    \smallskip
    Il null split viene selezionato solo se \texttt{len(candidates)~<~2}
    (unico candidato). Con \texttt{len~$\ge$~2}, il controllo
    \texttt{best.merit~>~0} fallisce prima ancora di raggiungere
    la selezione del candidato — il null non viene mai scelto.
  \end{block}

  \medskip
  \begin{alertblock}{Quando cambia davvero il comportamento}
    \small
    Con feature continue il null split non può diventare il migliore e
    l'opzione è normalmente ininfluente. Con una feature nominale ad alta
    cardinalità, invece, i branch possono avere merit zero:
    \begin{itemize}
      \item senza pre-pruning: split forzato anche senza segnale;
      \item con pre-pruning: il null split disattiva la foglia ed evita
            lo split spurio.
    \end{itemize}
    \textbf{Test mirato:} 60 foglie senza pre-pruning, 1 foglia con pre-pruning;
    River e CapyMOA coincidono a ogni campione.
  \end{alertblock}
\end{frame}

\begin{frame}{La classe principale orchestra quattro responsabilità}
  \begin{columns}[T]
    \column{0.48\linewidth}
    \begin{block}{Configurazione}
      \small
      Traduce le opzioni MOA in parametri operativi e crea i prototipi
      di observer e drift detector.
    \end{block}

    \begin{block}{Training}
      \small
      Inizializza la radice e delega ogni istanza a
      \texttt{AdaLeafNode} oppure \texttt{AdaBranchNode}.
    \end{block}

    \column{0.48\linewidth}
    \begin{block}{Predizione}
      \small
      Raccoglie le foglie raggiunte nel path principale e negli alberi
      alternativi, quindi ne media le predizioni.
    \end{block}

    \begin{alertblock}{Crescita e memoria}
      \small
      Decide quando creare uno split, aggiorna i contatori e applica
      periodicamente il limite di memoria.
    \end{alertblock}
  \end{columns}
\end{frame}

\begin{frame}{Memory management}
  \begin{columns}[T]
    \column{0.48\linewidth}
    Ogni $10^6$ istanze (configurabile), la classe principale esegue un
    ciclo di stima e controllo della memoria:

    \medskip
    \begin{enumerate}
      \item \textbf{Stima} (\texttt{estimateModelByteSizes}):
            misura la RAM reale via \texttt{SizeOf} e aggiorna le stime
            per foglia attiva e inattiva
      \item \textbf{Controllo} (\texttt{enforceTrackerLimit}):
            se l'albero supera \texttt{maxByteSize}, ordina le foglie per
            \emph{promise} (profondità) e \textbf{disattiva} le meno
            promettenti finché si rientra nel limite
      \item \textbf{Potatura observer} (\texttt{pruneLeafObservers}):
            dopo un tentativo di split non riuscito, rimuove dagli EBST
            i tagli ormai troppo deboli rispetto al migliore
    \end{enumerate}

    \column{0.48\linewidth}
    \begin{alertblock}{Foglia attiva vs.\ inattiva}
      Una foglia \textbf{inattiva} smette di aggiornare i propri
      attribute observer (risparmio di RAM), ma continua ad
      aggiornare statistiche e modello di predizione.
      Non può più splittare.\\[4pt]
      Se il budget lo permette, viene \textbf{riattivata}.
    \end{alertblock}

    \medskip
    \begin{block}{Opzione \texttt{stopMemManagement}}
      Se attiva, al superamento del limite si blocca direttamente
      \texttt{growthAllowed} invece di deattivare foglie: nessun
      nuovo split viene tentato.
    \end{block}

    \medskip
    \textcolor{gray}{\small \texttt{SizeOf} richiede l'agente JVM
    \texttt{-javaagent:sizeofag.jar}; se assente, il meccanismo è
    silenziosamente disabilitato.}
  \end{columns}
\end{frame}

\begin{frame}{Dal drift allo swap: il ciclo in cinque passi}
  \begin{columns}[T]
    \column{0.56\linewidth}
    \begin{enumerate}
      \item \textbf{Misura l'errore}\
            Predizione del path principale e aggiornamento di
            $|y-\hat y|$.
      \item \textbf{Interroga ADWIN}\
            Una variazione significativa segnala un possibile drift.
      \item \textbf{Crea l'alternativa}\
            Una nuova foglia nasce alla stessa profondità e apprende
            in parallelo.
      \item \textbf{Attende evidenza}\
            Il confronto parte solo quando entrambi gli error tracker
            superano la soglia minima.
      \item \textbf{Decide}\
            Lo z-test produce uno \emph{swap} oppure elimina l'alternativa.
    \end{enumerate}

    \column{0.40\linewidth}
    \begin{alertblock}{Ordine degli aggiornamenti}
      \small
      L'istanza viene propagata sia all'albero alternativo sia al figlio
      corrente. I due modelli osservano quindi lo stesso stream dal momento
      in cui l'alternativa viene creata.
    \end{alertblock}

    \medskip
  \end{columns}
\end{frame}

\begin{frame}[fragile,shrink=5]{Codice: nascita e selezione dell'albero alternativo}
  \begin{columns}[T]
    \column{0.48\linewidth}
    \textbf{1. Rilevazione e creazione}
\begin{lstlisting}[style=java,basicstyle=\ttfamily\fontsize{6.9}{8.2}\selectfont]
double err = Math.abs(y - yPred);
double oldMean = errorTracker.getMean();

driftDetector.update(err);
errorTracker.update(err, 1.0);
boolean drift = driftDetector.isDrift();

// Ignore changes while error improves
if (drift && errorTracker.getMean() < oldMean) {
  errorTracker = new VarStats();
  drift = false;
}

if (drift && alternateTree == null) {
  errorTracker = new VarStats();
  alternateTree = tree.newLeaf(null, depth);
  tree.nAlternateTrees++;
}
\end{lstlisting}

    \column{0.48\linewidth}
    \textbf{2. Z-test e decisione}
\begin{lstlisting}[style=java,basicstyle=\ttfamily\fontsize{6.9}{8.2}\selectfont]
double denom = Math.sqrt(
    altV / altN + curV / curN);
double z = (denom > 1e-12)
    ? (altMu - curMu) / denom : 0.0;
double p = 2.0 * NormalDist.cdf(-Math.abs(z));

if (p <= tree.switchSignificance) {
  if (altMu < curMu) {
    // Alternate wins: replace current subtree
    parent.children.set(parentBranch,
                        alternateTree);
    tree.nSwitchAlternateTrees++;
  } else {
    // Current subtree wins: discard alternate
    killNode(alternateTree, tree);
    alternateTree = null;
    tree.nPrunedAlternateTrees++;
  }
}
\end{lstlisting}
  \end{columns}

  \smallskip
  \centering
  \footnotesize
  Il confronto viene eseguito solo quando entrambi gli error tracker
  superano \texttt{driftWindowThreshold}.
\end{frame}

\begin{frame}{Compatibilità API: River $\to$ CapyMOA HATR}
  \vspace{-0.3em}
  \footnotesize
  \begin{columns}[T]
    \column{0.50\linewidth}
    \begin{tabular}{@{}ll@{\;}c@{}}
      \toprule
      \textbf{Parametro River} & \textbf{CapyMOA} & \\
      \midrule
      \texttt{grace\_period}            & \texttt{grace\_period}            & \textcolor{adwingreen}{$\checkmark$} \\
      \texttt{delta}                    & \texttt{split\_confidence}        & \textcolor{adwingreen}{$\checkmark$} \\
      \texttt{tau}                      & \texttt{tie\_threshold}           & \textcolor{adwingreen}{$\checkmark$} \\
      \texttt{leaf\_prediction}         & \texttt{leaf\_prediction}         & \textcolor{adwingreen}{$\checkmark$} \\
      \texttt{model\_selector\_decay}   & \texttt{model\_selector\_decay}   & \textcolor{adwingreen}{$\checkmark$} \\
      \texttt{bootstrap\_sampling}      & \texttt{bootstrap\_sampling}      & \textcolor{adwingreen}{$\checkmark$} \\
      \texttt{min\_samples\_split}      & \texttt{min\_samples\_split}      & \textcolor{adwingreen}{$\checkmark$} \\
      \texttt{drift\_window\_threshold} & \texttt{drift\_window\_threshold} & \textcolor{adwingreen}{$\checkmark$} \\
      \texttt{switch\_significance}     & \texttt{switch\_significance}     & \textcolor{adwingreen}{$\checkmark$} \\
      \texttt{max\_depth}               & \texttt{max\_depth}               & \textcolor{adwingreen}{$\checkmark$} \\
      \texttt{merit\_preprune}          & \texttt{merit\_preprune}          & \textcolor{adwingreen}{$\checkmark$} \\
      \texttt{remove\_poor\_attrs}      & \texttt{remove\_poor\_attrs}      & \textcolor{adwingreen}{$\checkmark$} \\
      \texttt{stop\_mem\_management}    & \texttt{stop\_mem\_management}    & \textcolor{adwingreen}{$\checkmark$} \\
      \texttt{seed}                     & \texttt{random\_seed}             & \textcolor{adwingreen}{$\checkmark$} \\
      \bottomrule
    \end{tabular}

    \column{0.47\linewidth}
    \begin{tabular}{@{}ll@{\;}c@{}}
      \toprule
      \textbf{Parametro River} & \textbf{Stato} & \\
      \midrule
      \multicolumn{3}{@{}l@{}}{\textit{Hardcoded al default River:}} \\[1pt]
      \texttt{splitter}        & fisso \texttt{TEBST}      & \textcolor{capyorange}{$\sim$} \\
      \texttt{drift\_detector} & fisso \texttt{ADWIN}      & \textcolor{capyorange}{$\sim$} \\
      \texttt{leaf\_model}     & fisso \texttt{Perceptron} & \textcolor{capyorange}{$\sim$} \\
      \midrule
      \multicolumn{3}{@{}l@{}}{\textit{Gestito automaticamente:}} \\[1pt]
      \texttt{nominal\_attrs}  & schema ARFF               & \textcolor{adwingreen}{$\checkmark$} \\
      \midrule
      \multicolumn{3}{@{}l@{}}{\textit{Non implementato:}} \\[1pt]
      \texttt{binary\_split}   & solo \texttt{False}       & \textcolor{driftred}{$\times$} \\
      \bottomrule
    \end{tabular}

    \medskip
    \textcolor{adwingreen}{$\checkmark$}~Opzioni principali esposte nel wrapper.\\
    \textcolor{capyorange}{$\sim$}~Tre componenti fissati al default River.\\
    \textcolor{driftred}{$\times$}~Una modalità non-default non implementata.

    \medskip
    \begin{alertblock}{}
      \scriptsize
      La configurazione standard è replicata. L'equivalenza puntuale vale
      senza bootstrap; con bootstrap si confrontano le metriche aggregate.
    \end{alertblock}
  \end{columns}
\end{frame}

% ============================================================
\section{Conclusioni}
% ============================================================

\begin{frame}{Cosa porta questo lavoro}
  \begin{columns}[T]
    \column{0.48\linewidth}
    \begin{block}{1. Un learner nativo}
      \small
      HATR entra nell'ecosistema MOA/CapyMOA senza delegare il training
      a River.
    \end{block}

    \begin{block}{2. Fedeltà verificabile}
      \small
      Nella configurazione deterministica, struttura e predizioni
      coincidono con River lungo tutto lo stream.
    \end{block}

    \column{0.48\linewidth}
    \begin{block}{3. Beneficio prestazionale}
      \small
      Nel benchmark il tempo complessivo migliora di circa
      2.95--7.85$\times$, soprattutto grazie al training nella JVM.
    \end{block}

    \begin{alertblock}{4. Un confine chiaro}
      \small
      Il port replica il percorso standard, ma non rende ancora
      pluggabili splitter, drift detector e modello di foglia.
    \end{alertblock}
  \end{columns}

  \medskip
  \centering
  \large
  \textbf{Risultato: HATR diventa confrontabile con gli altri learner
  streaming direttamente in CapyMOA.}
\end{frame}

\begin{frame}{Sviluppi futuri}
  \begin{columns}[T]
    \column{0.50\linewidth}
    \begin{block}{Completamento parametri non-default}
      \begin{enumerate}\small
        \item \textbf{\texttt{binary\_split=True}} \\
              \footnotesize Modificare \texttt{HANominalObserver}
              per imporre split binari anche su attributi nominali;
              aggiungere flag \texttt{-B} al learner MOA e
              parametro nel wrapper Python.

        \medskip
        \item \textbf{Splitter pluggabile + \texttt{AdaNumMultiwayBranch}} \\
              \footnotesize Implementare \texttt{AdaNumMultiwayBranchReg}
              e un \texttt{HAQOObserver} per supportare
              \texttt{QOSplitter(allow\_multiway\_splits=True)};
              esporre \texttt{splitter} come opzione MOA.

        \medskip
        \item \textbf{Leaf model pluggabile} \\
              \footnotesize Aggiungere supporto per modelli leaf
              alternativi al \texttt{Perceptron}
              (es.\ regressione ridge).
      \end{enumerate}
    \end{block}

    \column{0.47\linewidth}
    \begin{block}{Estensioni di architettura}
      \begin{enumerate}\setcounter{enumi}{3}\small
        \item \textbf{Drift detector pluggabile} \\
              \footnotesize Definire un'interfaccia
              \texttt{HADriftDetector} e supportare detector
              alternativi ad \adwin{}
              (es.\ \textsc{eddm}, Page-Hinkley).
      \end{enumerate}
    \end{block}
  \end{columns}
\end{frame}

\begin{frame}[allowframebreaks]{Riferimenti bibliografici}
  \begin{thebibliography}{9}
    \setbeamertemplate{bibliography item}[text]

    \bibitem{Domingos2000}
    P.~Domingos and G.~Hulten.
    \newblock Mining high-speed data streams.
    \newblock \emph{KDD}, 2000.

    \bibitem{Bifet2007}
    A.~Bifet and R.~Gavaldà.
    \newblock Learning from time-changing data with adaptive windowing.
    \newblock \emph{SDM}, 2007.

    \bibitem{Bifet2009}
    A.~Bifet and R.~Gavaldà.
    \newblock Adaptive learning from evolving data streams.
    \newblock \emph{IDA}, 2009.

    \bibitem{Babcock2003}
    B.~Babcock, M.~Datar, R.~Motwani, and L.~O'Callaghan.
    \newblock Maintaining variance and k-medians over data stream windows.
    \newblock \emph{PODS}, 2003.

    \bibitem{Gomes2017}
    H.~M.~Gomes \emph{et al.}
    \newblock Adaptive random forests for evolving data stream classification.
    \newblock \emph{Machine Learning}, 2017.

    \bibitem{Ikonomovska2011}
    E.~Ikonomovska, J.~Gama, and S.~Dzeroski.
    \newblock Learning model trees from evolving data streams.
    \newblock \emph{Data Mining and Knowledge Discovery}, 2011.

    \bibitem{MOA2010}
    A.~Bifet \emph{et al.}
    \newblock MOA: Massive Online Analysis.
    \newblock \emph{JMLR}, 2010.

    \bibitem{CapyMOA2024}
    C.~Grzenda \emph{et al.}
    \newblock CapyMOA: A Python library for data stream mining.
    \newblock \emph{arXiv}, 2024.
  \end{thebibliography}
\end{frame}

\begin{frame}
  \begin{center}
    \Large\textbf{Grazie per l'attenzione}

    \bigskip
    \normalsize
    Codice, notebook e JAR disponibili in:\\
    \texttt{2025-2026\_Regression-Models-in-CapyMOA/andrea\_zhang/}

    \bigskip
    \begin{tabular}{ll}
      Implementazione: & \texttt{implementation/hatr-moa/} \\
      Wrapper Python:  & \texttt{implementation/hatr\_capymoa/} \\
      Esperimenti:     & \texttt{experiments/} \\
      Tutorial:        & \texttt{tutorials/} \\
    \end{tabular}
  \end{center}
\end{frame}

\end{document}
